\chapter{Why Do We Keep Asking ``Who Are We?''}
\section{An Old Photograph at the Bottom of a Drawer}
First, pick something up. Not a bronze artifact in a museum or a famous book in a library, but something many families have at home: an old photograph.

Imagine this scene. Your grandmother is moving, and you are helping clear out her aging chest of drawers. One drawer sticks. You give it a hard pull, and a yellowing photograph slips out from underneath and floats to the floor. You pick it up and blow off the dust.

The photograph is black and white, its four corners worn round. Seven or eight people stand in front of a brick wall. The faces in the back row are a little blurred. In front, a small boy holds an enamel cup with what looks like writing on it, though you cannot make it out. At the far left stands a young man in a white shirt, his pose rather stiff, as though he is not used to facing a camera. On the back is a line written in fountain pen. The ink has bled a little, but you can just decipher ``Summer 1974,'' followed by a place name.

Your grandmother leans over for a look. ``Oh, we still have this one!'' she says.

Now pay attention to what happens next, because it is the starting point of this chapter and of the whole book. Holding this photograph, you can ask many questions, and those questions sort themselves into several piles. One pile the photograph answers immediately: How many people are in it? Seven or eight? What are they wearing? What building stands behind them? Those answers are printed on the photographic paper. All you need is a pair of eyes. Another pile the photograph can answer with the help of a cup of tea and an afternoon with Grandma: What year was this? How were these people related? Why had they gathered that day?

But about another pile of questions, the photograph stays silent. Who stood behind the camera? Why was the young man on the far left smiling so reluctantly? What had the group been talking about five minutes before the picture was taken? And if you look closely, a narrow strip has been cropped from the right-hand edge. Was someone else originally standing there?

The deepest question comes last. You ask your grandmother, ``Who are they?'' She says, ``That's our family.'' But ``our family'' does not close the question. It opens it: What does ``us'' mean? What binds these seven people, with their different expressions and their utterly different later lives, into an ``us''? And why should you, a child born decades later, feel a little stir inside on seeing the photograph, a sense that it has something to do with you?

A photograph can tell us a great deal, but it hides even more. It pins some things to paper while letting others drift away forever. The humanities, the whole field of knowledge this book explores, grow out of that feeling of wanting to know more.

\section{Rain Falls on a Small Town}
Now come with me into a scene. This time, instead of a photograph, there is a rainstorm.

It is a June afternoon in a riverside town in southern China. The morning was sultry, the air as heavy as a damp quilt, the cicadas calling listlessly. After noon, clouds began piling up beyond the river, turning from white to gray and from gray to a gleaming leaden black. The wind suddenly changed direction, carrying the smell of water. Street vendors began gathering up their plastic sheets. Then the first raindrop struck a stone paving slab, spreading into a dark circle. Within ten seconds, countless circles merged. The rain had arrived.

The same rain, in the same minute, falls on four people in the town.

The first is at the weather station on the edge of town. She watches the radar image on her screen, where the red patch representing rainfall is moving slowly. She records the time and intensity, estimates the amount of rain this weather event will bring, and sends the figures out. For her, the rain is a \textbf{meteorological process}: warm, moist air, convection, precipitation, a chain of physical events that can be measured and predicted. When she says, ``Thirty millimeters of rain fell,'' instruments anywhere in the world can check whether she is right.

The second person is in a rice field across the river. He looks up at the rain, catches some in his hand, then bends to squeeze the soil. There has been a drought for more than twenty days, and the tips of the rice seedlings have begun to yellow. For him, this rain is not data. It is \textbf{a lifeline}: his child's school fees for the second half of the year, the date a loan comes due, whether the family can breathe a little easier around the dinner table. Or he may frown: if the rain falls too hard and too fast, the newly transplanted seedlings will be ruined. For him, this rain is \textbf{a farming matter}, a livelihood, something demanding an immediate decision: should he go out tonight to check the water in the field?

The third person stands outside the town's temple. Last month, when the drought was severe, the town's elders held a rainmaking ritual here. Now the rain has come. Someone bows toward the temple door and says, ``Our prayers worked.'' For him, the rain is \textbf{an answer}, an exchange between human beings and some greater power. What it proves is not the existence of air currents, but that ``we have not been forgotten.''

The fourth is a student at the town's secondary school, daydreaming in a window seat. Raindrops tap against the glass. The sports field quickly turns gray-white, and the curtain of rain slowly erases the distant hills. Something suddenly stirs inside her. She takes out a notebook and writes a sentence about rain in a corner. It is not very good, and one day she will tear it out, but right now it hurts not to write it. For her, the rain is \textbf{a poetic image}, something she cannot quite explain and can only hold in a poem.

It is still the same cloud and the same water molecules, falling from the same sky onto the same town. Yet the rain is four things at once: a meteorological process, a farming matter, an answered prayer, a poetic image. Which is its ``real'' identity?

\section{Which Is the Truth About the Rain?}
Do not answer just yet. Let us make the conflict a little clearer.

One powerful argument says that the weather station's account is the \textbf{truth} about the rain. The other three are merely things people add to it: the farmer's anxiety, the elder's superstition, the student's sentimentality. The rain is ``actually'' just water vapor condensing and falling, thirty millimeters of it, nothing more. Everything else is projection, story, mood.

That sounds brisk and scientific, but it has a subtle flaw: it quietly replaces the word ``truth'' with ``physical description.'' Notice that when the farmer says, ``This rain saved my rice,'' his statement \textbf{is not false}. The rice really suffered from drought, the rain really came, and his livelihood really changed as a result. When the elder says, ``Our prayers worked,'' you may disagree with his explanatory framework, but the sequence he describes, ritual, drought, and rain, really happened. As for the student's line of poetry, true or false is beside the point: it was never meant to be a weather report. These four accounts are not competing in the same race, so none simply eliminates the others.

The opposite mistake is just as real. If the temple elder declares, ``Our prayers brought the rain; the weather station's figures don't matter,'' the figures will quickly teach him otherwise when tomorrow he must decide whether to dig irrigation channels or drain the fields. If the student insists that ``the meaning of rain exists only in the heart,'' she will still have to bring in the washing from the balcony when she gets home.

The real question now emerges. It is not ``Which account is correct?'' but \textbf{Why does the same event branch into several different kinds of question? How are those questions related?}

At this point, we can sketch the territory of the humanistic knowledge this book explores. Broadly speaking, it studies four things:

\begin{itemize}
\item \textbf{Meaning}: what things mean to people. Rain is a livelihood to the farmer, an answer to the elder, a nuisance to you passing with an umbrella. The meaning of the same event varies with the person and the situation.
\item \textbf{Action}: what people do with purposes and reasons. Whether the farmer checks the water tonight is not molecular motion; it is a decision.
\item \textbf{Institutions}: arrangements people establish together that then constrain them. Temples, schools, timetables, marriage, money, states: when a bell rings and the whole class stands up, that is not a law of physics. It is a convention everyone tacitly agrees to follow.
\item \textbf{Expression}: things people make to give meaning a form. A line of poetry, a photograph, a ritual, a song.
\end{itemize}
Meteorological processes belong to the natural sciences, not the humanities. But these two fields are not rivals. They divide the work and cooperate. A weather forecast helps the farmer decide what to do: that is cooperation. A doctor uses instruments to treat illness, but treatment may be less effective without understanding how patients see their illness: whether they are frightened, whether they trust the treatment, whether they are hiding the illness from their families. That, too, is cooperation. We will see repeatedly that hard questions often need both sides working together.

Why do humanistic questions form a continent of their own? The reason is simple: \textbf{rain does not ask who it is, but people do.} Rain does not care whether you call it a meteorological process or an answered prayer. It falls all the same. People are different. They explain their actions, tell others those explanations, and change what they do next because of an explanation. When you study rain, it neither cooperates nor protests. When you study people, they talk back. That is the entrance to the question ``Who are we?'' and the place this opening chapter really wants to explore.

\section{Four Kinds of ``Us,'' Four Answers}
Human beings have asked ``Who are we?'' for thousands of years, and their answers differ widely. Let us first listen to voices from around the world, then return to the hardest dispute.

In New Zealand, Māori people have a traditional way of introducing themselves called a pepeha. A Māori person introducing themselves does not begin with a name and occupation, but with their mountain, their river, the canoe on which their ancestors reached this land, and the branch to which their family belongs. Only after the mountain, river, canoe, and ancestors comes their own name. In this answer, ``Who am I?'' is first a question about \textbf{where I come from}. I am not an isolated individual. I am a link in a long chain whose other end is fastened to the land and my ancestors.

In ancient Greece, a famous maxim was inscribed on Apollo's temple at Delphi: ``Know yourself.'' This points in another direction. Instead of supplying an answer, it defines a person as \textbf{a puzzle requiring self-examination}. You do not first establish who you are and then begin living. You must spend a lifetime asking, slowly finding out. Many figures in Greek tragedy come to grief precisely because they think they already know who they are.

In southern Africa, an influential idea known in Zulu as ubuntu means, roughly, that a person becomes a person through other people. It is often paraphrased as ``I am because we are.'' In this understanding, relationships are not additions outside a person. They \textbf{are} part of what constitutes that person. Cut all your connections to others, and what remains is not a ``pure self'' but an incomplete human being.

In Chinese tradition, the answer may sound familiar. Confucianism never considers a person in isolation, but places them in a network of relationships: someone's child, someone's parent, someone's friend, someone's student. Remove the network, and ``Who are you?'' becomes unanswerable. This echoes Māori traditions and ubuntu from afar, while growing in soil of its own.

Put these four answers together and two things emerge. First, ``Who are we?'' has never had one standard answer. Each answer is a way of living shaped by a tradition over thousands of years. Second, the answer you may find most self-evident, ``I am an independent individual; first and foremost, I am myself,'' is actually one among many, and a very young one. Treating it as humanity's default setting is itself a prejudice.

But a question follows: given so many different answers, \textbf{how should we study} the humans who give them? Here are two genuinely opposing approaches. Let me make the strongest case for each, especially the one that often receives less sympathy.

\textbf{Position one: We should study people as we study nature, through measurement and experiment.}

This position is often mocked today as ``cold'' or ``inhuman.'' That is unfair; its reasons are substantial. First, measurement is \textbf{testable}. You say a teaching method works; I say it does not. Argument alone may never settle it. Compare groups drawn from several thousand students, and the answer can emerge. Second, measurement \textbf{saves lives}. Whether a vaccine works or a medicine is effective cannot be established by ``understanding patients' feelings.'' It requires rigorous controlled experiments. The eras when illness was treated as a moral failing or divine punishment were precisely those with the highest human death rates. Third, making people ``special'' has sometimes obstructed progress. Once people are declared sacred and beyond investigation, anatomy is prohibited, psychology is kept outside, and much suffering is dismissed with ``It was meant to be.'' What makes the scientific method valuable is precisely that it spares no authority, including the idea of human exceptionalism.

\textbf{Position two: Measurement alone loses the crucial fact that people interpret their own actions.}

This position begins with a simple observation. Consider the same movement: someone raises a hand. It might be a greeting, a request to speak in class, an attempt to hail a taxi, a vote, or a shoulder cramp. The muscular movements are almost identical, but what determines what happens next is their \textbf{meaning}. Whether the driver stops or the teacher calls on the student depends on what the movement ``counts as.'' Meaning is not in the muscles. It is in the setting, the rules, and the shared understanding of those involved. A high-speed camera can record every detail of a wave, but it can never film ``This is a vote in favor.''

There is a deeper point. People are unlike raindrops. When you study raindrops' motion, they do not read your paper and change their trajectories. But publish a report saying that ``young people today are increasingly unwilling to marry,'' and young people will read it and argue about it. Some may become angry and deliberately do the opposite; others may use it to persuade their parents. \textbf{The research itself enters what it studies, and its subjects talk back.} This has no counterpart in the natural sciences, yet humanistic research faces it every day.

Both positions have genuine reasons on their side. This book does not choose one and discard the other. It asks which tools a particular kind of question requires. That brings us to this chapter's tool.

\section{Thinking Bridge: What Kind of Question?}
Faced with a rainstorm, a photograph, or any argument, here is a tool you can take away immediately: \textbf{Before speaking, first ask what kind of question this is.}

The questions people ask about an event fall broadly into four categories.

\textbf{First: questions of fact. What happened?} How many people are in the photograph? How long did the rain last? On what date was the rainmaking ritual held? In principle, such questions have an answer that holds for everyone, and evidence, from eyes, instruments, archives, or records, decides it. A factual question may be extremely hard or may never be settled, but there \textbf{is} such a thing as a correct answer.

\textbf{Second: questions of explanation. Why did it happen?} Why did the whole family gather that year? Why did this town's rainmaking ritual stop once the prayers were over? These questions concern causes. There is often more than one candidate answer, and we must compare how well the explanations hold up.

\textbf{Third: questions of meaning. What did it mean to those involved?} What did the rain mean to the farmer, the elder, and the student? No instrument can do this work for you. You can approach an answer only through their words, actions, and creations. Notice that answers about meaning \textbf{can also be right or wrong}. You cannot arbitrarily announce that the farmer regarded the rain as an enemy; evidence would contradict you. But this evidence is interpreted through understanding, not obtained through measurement.

\textbf{Fourth: questions of evaluation. How should we judge it?} Should the rainmaking ritual be mocked? Is including ancestors in a self-introduction a good thing or a constraint? Ultimately, such questions require value judgments, and value judgments need \textbf{reasons}, not merely a loud voice.

Use this tool to ``sort'' an argument. Many quarrels never end because the two sides are not answering the same kind of question. One answers a factual question: ``The rainfall was thirty millimeters.'' The other answers a question of meaning: ``You have no idea what this rain means to me.'' Both statements may be true, yet they are treated as contradictions. Once sorted, most online bickering either dissolves or becomes the question actually worth discussing.

Be especially alert when someone \textbf{deliberately} confuses the categories. When someone says, ``Young people just can't endure hardship,'' in the tone of a factual report, they have quietly smuggled in an evaluation. When an advertisement describes a drink in poetic language, it disguises meaning as fact. Learn to sort the categories, and you equip yourself with an all-weather radar against being fooled.

This tool runs through the entire book. In every later chapter we will repeatedly distinguish what happened (fact), why it happened (explanation), how people at the time understood it (meaning), and how we judge it (value). Now let us use it to read a short primary source from two thousand years ago.

\section{On the Hao River Bridge: ``How Do You Know?''}
During China's Warring States period, Zhuangzi and his friend Huizi had a famous debate, recorded in the ``Autumn Floods'' chapter of the Zhuangzi. They were walking on a bridge over the Hao River:

\begin{quote}
Zhuangzi said, ``The minnows swim out at ease. Such is the joy of fish.'' Huizi said, ``You are not a fish. How do you know the joy of fish?'' Zhuangzi said, ``You are not I. How do you know that I do not know the joy of fish?''
\end{quote}
In everyday language, Zhuangzi says: Look how leisurely that fish swims. It is happy. Huizi immediately replies: You are not a fish. How do you know it is happy? Zhuangzi is just as quick: You are not me. How do you know I do not know it is happy?

On the surface, this looks like two intellectuals with nothing better to do than bicker. Apply our sorting tool, however, and the debate becomes remarkably sharp. Huizi's attack targets a vital point in humanistic inquiry: \textbf{What entitles us to know what is inside another mind?} You are not a fish. Your ``The fish is happy'' is conjecture, projection, your own mood pasted onto the fish. You might similarly say, ``The person in this photograph looks happy.'' But how do you know? Perhaps they had just quarreled with their family when the picture was taken.

Huizi's question is both severe and well aimed. If it holds, every statement about another person's inner life, ``He is sad,'' ``She is sincerely apologizing,'' ``This family seems harmonious,'' loses its footing. The humanities might as well close up shop.

The cleverness of Zhuangzi's reply is that he does not meet the attack head-on. He \textbf{turns its skepticism back}. You say that because I am not a fish, I cannot know a fish. But you are not me, so how can you be certain that I do not know? You attack me with doubt, yet doubt itself must first cross the barrier of knowing another's mind. To question me, you must first understand what I am saying, and understanding speech is already a kind of crossing between minds.

This debate has gone undecided for more than two thousand years, and rightly so, because each side grasps half a truth. Huizi's half is that \textbf{understanding another can never be as direct as measuring temperature}. Inference always stands between us, and we must be honest about that distance. This is what our opening photograph teaches: the brighter its subjects smile, the less entitled we are to announce that we ``know'' they are happy. Zhuangzi's half is that \textbf{if Huizi's skepticism held completely, all understanding between people, including this debate itself, would be impossible}. Yet the debate plainly occurred, and every day we genuinely understand one another to a considerable degree. That understanding is indirect and fallible, but it really works.

Two thousand years later, the same argument continues with new faces. Can an animal feel pain? What exactly did the author of a text intend? Do the words spoken by artificial intelligence count as ``understanding''? Whenever you encounter such disputes, remember the two people on the Hao River bridge. Some of humanity's sharpest minds have been pulling back and forth here for a very long time.

\section{Why Do People Keep Telling Stories About ``Us''?}
Return now to a fact that evidence can support: almost every known human society has left stories about who ``we'' are and where ``we'' came from. Genealogies, creation myths, heroic epics, founding legends, school histories, factory histories, dinner-table accounts of ``what your grandfather did back then'': the forms vary enormously, but the practice is widespread. Nobody composes a set of weather data to explain ``where we came from.'' A society without ``our story'' is like a person without memory.

The fact is clear, but \textbf{why} is it so? At least three genuinely different explanations stand before us. They are not interchangeable; each has its reasons and its limits. Let us set them out one by one. Remember, we are comparing answers to ``Why does this happen?'' not to ``Are these stories true?''

\textbf{Explanation one: This is how the brain works. Humans are storytelling machines.}

This explanation starts with individual psychology. Observe yourself. When recalling yesterday, you remember not a chronological inventory of events, but small stories with causes and turns: ``I nearly arrived late because my alarm didn't ring, but luckily I made it.'' The human brain naturally organizes fragments into plots with beginnings and endings. Without that ability, even ``Who am I?'' becomes hard to sustain. For someone who loses their memory, the greatest pain is not lost information but a lost personal story. This explanation has the advantage of being testable in every person, and it explains why the urge for stories appears so naturally and so early. Soon after learning to speak, children pester adults to ``tell it again.'' Its limit is also clear: the brain can explain why \textbf{individuals} love stories, but not why stories are \textbf{collective}. Why does a whole people hold on to one story for thousands of years? The mechanism of individual memory does not reach the scale of collective identity.

\textbf{Explanation two: Stories are the glue of cooperation. Shared stories turn strangers into an ``us.''}

This explanation starts with society. A person can know only a limited number of people face to face, yet humans organize cities of tens of thousands and countries of hundreds of millions. How? By believing in the same invisible things: we descend from one ancestor, are citizens of one country, attended one school. These shared stories act as glue, joining strangers into groups that can farm, fight, and pay taxes together. This is a powerful explanation. It shows why societies expend so much effort on stories: writing histories, raising monuments, celebrating festivals, commemorating events. Its limit is that it makes stories look too much like \textbf{tools}. If stories are merely glue, why are people willing to weep for them or die for them? Glue does not make someone sit late at night gazing at an old photograph. A purely practical balance sheet cannot account for emotion.

\textbf{Explanation three: Stories mark power's boundaries. Defining ``us'' also defines ``them.''}

This explanation changes our vantage point. It notices that every ``us'' excludes someone. Saying ``We are family'' also says that others are not; saying ``We are insiders'' marks outsiders. So \textbf{who is entitled to tell the story, and whose version counts}, becomes crucial. Who enters a family's genealogy, and who is quietly omitted? Who occupies the center of a nation's narrative, and who appears only in a footnote? None of this is neutral. This explanation is sharp enough to illuminate corners the other two miss. But it has blind spots of its own. Reduce every story to an exercise of power, and you cannot explain the genuine comfort and belonging stories contain. The smile of the child holding the enamel cup was not entirely arranged by someone else.

All three explanations grasp part of the truth, and none reaches all of it. Remember that feeling; it will run through this book. When an explanation claims to explain everything, it has often discarded something in advance.

\section{Six Lenses: Where This Book Goes Next}
After all this discussion of ``Who are we?'' we can now make clear how this book intends to tell its story.

The next twenty-odd chapters follow a timeline through the human story from beginning to end: early humans who could speak and tell stories; farming and cities; writing and states; the great classical civilizations; worldwide religious and philosophical traditions; exchanges across continents and the age of ocean voyages; science, revolution, and industrialization; and finally today's globalization, digital age, and artificial intelligence. This is the book's warp.

Its weft consists of six lenses. We will take turns viewing the same human phenomenon through each:

\begin{itemize}
\item \textbf{Language}: why people need language, and how it shapes thought and society. Through this lens, our rainstorm raises another question: How were words such as ``rain'' and ``us'' invented, and how do they in turn frame us?
\item \textbf{Literature}: why people tell stories, and why fiction can tell the truth. Through this lens, the student's line of poetry about the rain feels more convincing to her than any rainfall data. Why?
\item \textbf{Religion}: how people face death, suffering, and what lies beyond their control. Through this lens, the temple ritual is not a ``mistaken weather forecast,'' but a way people find their bearings amid immense uncertainty.
\item \textbf{Philosophy}: what knowledge deserves belief, and how to distinguish argument, evidence, and authority. Through this lens, the argument on the Hao River bridge has no loser; it traces the boundary between ``knowing'' and ``believing.''
\item \textbf{History}: how history happened, and what grounds our knowledge of the past. Through this lens, ``Summer 1974'' on the back of the photograph is a historical source. Judging its reliability requires a whole set of methods.
\item \textbf{Aesthetics}: why people create art, and what makes something move us. Through this lens, why might a passing stranger stop for a long time before this old photograph?
\end{itemize}
These six lenses are not six unrelated subjects. They are \textbf{six doors into the same house}. Remember the rain: meteorological process, farming matter, answered prayer, and poetic image at once. There is no way to see it all in a single glance. You must enter one door at a time. What you see through each is real, and none is the whole. A seemingly small question, such as ``Why does a joke hurt?'' can open six layers: the tone of words, the act of speaking, a group's power, literary irony, philosophical inquiry, and more. This book provides steps into each layer. Where those steps ultimately lead, the six lenses will take turns showing you.

In each chapter ahead, you will start from a particular object, a coin, a temple, a poem, a phone, pass through a scene, a dispute, tools, and evidence, and finally return to your life today. The old photograph will recede into the background, but its lesson will not expire: \textbf{Seeing a thing is easy; knowing what it means is hard. Yet people live precisely within that ``what it means.''}

\section{Deeper: Explanation and Understanding}
Now push open this chapter's deepest door. We will formally meet a distinction that helped lay the foundations for how the modern humanities understand themselves.

In late nineteenth-century Germany, the philosopher Wilhelm Dilthey spent his life considering a question: What exactly separates the study of people from the study of nature? His answer was, roughly, \textbf{Nature we explain; the human mind we understand.} These correspond to two ways of knowing. ``Explanation'' seeks regularities from outside. Water vaporizes when heated to its boiling point. You do not need to know what water ``thinks,'' because it does not think. You need only identify a regularity that holds everywhere. ``Understanding'' enters from within. You want to know what an action means to its actor, what a work means to its creator and audience. You must temporarily enter another person's situation and see the world through their eyes.

Later, the sociologist Max Weber brought this idea into the study of social action. Broadly, his thought was that sociology should not merely count the movements people make, but understand their \textbf{subjective meaning} for the actors. Only by understanding meaning can we explain action. Return to the raised hand. A camera can give you the whole truth at the level of ``explanation,'' but the fact that ``this is a vote in favor'' exists only at the level of ``understanding.'' Remove that level, and most facts of human society, voting, weddings, classes, loan repayments, evaporate on the spot, leaving only bodily movements.

Notice an easy misjudgment; it is also this chapter's \textbf{counterexample}. After hearing understanding emphasized, you might conclude that people's explanations of themselves are the final answer: whatever someone says they mean is what they mean. \textbf{That inference is wrong.} Return to the old photograph to see why. If everyone in the family describes the day it was taken, you will probably hear conflicting versions, each offered sincerely and confidently. This need not mean anyone is lying. Human memory and self-explanation are unreliable in themselves. In the late twentieth century, a series of studies by the psychologist Elizabeth Loftus and others repeatedly showed that people can not only remember incorrectly, but under suggestion ``remember'' events that never happened and believe them firmly. Someone may be led to ``remember'' being lost in a shopping mall as a child, with vivid details and tears, though it never occurred. The point is that our inner narrator is not a tape recorder, but a screenwriter continually revising the script.

This counterexample matters enormously because it keeps humanistic inquiry from two kinds of naivety. The first is ``scientific naivety'': since self-explanation is unreliable, trust only instruments and discard meaning entirely. The second is ``romantic naivety'': since we must understand people, believe every word they say. Humanistic inquiry belongs between them: \textbf{Take people's self-explanations seriously without accepting them wholesale.} Listen to what the farmer says the rain means, but also check the harvest. Listen to why someone says they acted, but also examine what they actually did. Meaning requires understanding; facts require checking. Neither skill can be spared. That makes studying people harder than studying rain. It also makes the enterprise honest: it never pretends to possess an instrument that can measure everything.

\section{Today: Labels, Tests, and Answering Machines}
This question, thousands of years old, is buzzing in your pocket right now.

Look around. A fourteen-year-old today can encounter more answers to ``Who are we?'' than anyone in any earlier age, and most come packaged for easy collection. Take a test and the screen announces your personality type; four letters then hang in your profile. Fill out a questionnaire and a horoscope explains your temperament. Open a social platform and, from your likes, an algorithm silently answers ``Who are you?'' for you, then presents content that ``people like you'' enjoy. These labels share one feature: they compress the lifelong task demanded at the Delphic temple, knowing yourself, into three minutes.

Are all these labels useless? Not necessarily. Viewed with our sorting tool, some are clumsy attempts to answer \textbf{questions of meaning}. ``What kind of person am I?'' can offer belonging and ready-made words for emotions, a genuine comfort at a lonely age. But watch for them crossing the boundary and pretending to answer \textbf{questions of fact}. Four letters describe a self-image you are currently willing to claim, not a weather report about you. Be even more alert when you turn a label into destiny: ``This is just the kind of person I am, so I can't do that.'' You have slipped from knowing yourself into placing a ceiling on yourself. An adopted answer is, after all, someone else's gift.

Business and politics also understand the weight of ``Who are we?'' Brands no longer sell only drinks and sneakers; they sell ``what kind of person you are.'' Fan communities build a close-knit ``us'' from a shared love, while also marking out ``them.'' Arguments about where a country or people came from, and what sort of people they are, can still set the whole internet ablaze. Take the three explanations discussed earlier, stories as psychological needs, as glue, and as boundaries, and test them against each identity dispute you encounter today. See where each shines and where its light fails.

A new variable is entering the scene: machines are beginning to answer too. Artificial intelligence can write your self-introduction, imitate your voice in replies, and produce a fluent account when asked ``Who are you?'' The humanities' old boundary, ``People explain their own actions,'' has for the first time met a nonhuman object that also explains itself. Does a machine's explanation count as explanation? A later chapter will address that question specifically. For now, remember only this: when a counterpart that also speaks appears, \textbf{working out what a human being is becomes more urgent}, not more outdated.

Today's versions of the two-thousand-year-old Hao River debate play out daily: beneath a pet video, ``Is it really happy?''; beneath a trending story, ``Was that sincere, or just a public persona?''; beneath a review, ``Is this personality test reliable?'' This is not a field about the past. It is a field about \textbf{right now}. It just happens to have thousands of years of history.

\section{Over to You: What Does Not Asking Miss?}
Return to the photograph at the bottom of the drawer.

Your grandmother says, ``That's our family,'' and now you know how much rests behind those light words: a physical record instruments can check, a memory that may be mistaken, a story that binds strangers into a group, a boundary between ``us'' and ``them,'' and a question each person must spend a lifetime answering.

From beginning to end, this chapter has urged you to ask ``Who are we?'' Before closing the book, though, make a serious case for the other side, then answer this question:

\textbf{Can someone live a good life without ever asking ``Who am I? Who are we?''}

If you think so, give your reasons. The village farmer who never reads poetry, never looks at old photographs, and cares only about the water in his fields: is what he misses a necessity or a luxury? Could ``knowing yourself'' be an overrated modern obsession, a shackle made by people who think too much? After all, the fish never asks whether it is a fish. It swims at ease. Was that not precisely what Zhuangzi first noticed before saying, ``Such is the joy of fish''?

If you think not, give your reasons too. Is the ``us'' of someone who never asks really left empty? Or have they merely let others, tradition, algorithms, advertisements, surrounding voices, fill in the answer? Which is worse: an answer that is not your own, or no answer at all?

Notice that this question has no standard answer, and I have deliberately not chosen a side for you. Whichever side you take, however, something has quietly happened. The moment you considered the question seriously, you began doing what this book hopes to teach. Instead of collecting a ready-made ``us,'' you began making one yourself.

That is the entrance to the humanities. See you in the next chapter.
